\subsection{Testing}
\label{sec:impl_testing}

The backend uses \path{pytest}. Asynchronous tests are supported through \texttt{asyncio\_mode = "auto"}; the test root is \path{tests/}. The default suite measures branch coverage for the \path{abridgeai} package, requires at least 75\% coverage, and excludes tests marked \texttt{destructive} or \texttt{voice\_live}. This section documents the implemented test strategy and coverage boundaries rather than reporting an unsupported historical pass-count total.

\subsubsection{Native Interview Test Coverage}

The native interview runtime is tested at both deterministic and database-backed boundaries. Unit tests verify feature-flag routing, native-agent tool exposure, persistent conversational-state preparation, server-controlled state notes, answer grading, transcript persistence, typed-input processing, resume behavior, and hard-stop termination. Representative modules are \path{tests/unit/test\_interview\_native\_agent.py}, \path{tests/unit/test\_interview\_native\_scenarios.py}, \path{tests/unit/test\_interview\_native\_grading.py}, \path{tests/unit/test\_interview\_native\_transcript.py}, \path{tests/unit/test\_interview\_native\_text\_input.py}, and \path{tests/unit/test\_interview\_native\_resume\_and\_hard\_stop.py}.

These tests focus on the native authority boundary: the conversational LLM may generate phrasing, but server code retains question selection, runtime-state persistence, hint and follow-up limits, progression, and terminal submission. They also verify that typed turns use receipt-based idempotency while spoken recognition remains a best-effort media path rather than an exactly-once transport.

Database-backed integration tests exercise native seams that cannot be established by isolated unit tests. \path{tests/integration/test\_interview\_typed\_turn\_receipt.py} verifies receipt durability, idempotency, question linkage, and compare-and-set behavior; \path{tests/integration/test\_interview\_late\_fold\_reeval.py} verifies reevaluation when a receipt lands after an earlier verdict; and \path{tests/integration/test\_interview\_finish\_race.py} verifies concurrent learner finish and native hard-stop terminalization. Together, these tests validate the persistence and race-safety properties described in Section~\ref{sec:impl_interviews}.

\subsubsection{Continuous Integration and Live Voice Boundary}
\label{sec:impl_ci}

The normal test command is \texttt{uv run pytest}. In continuous integration, migrations are applied before one complete coverage run:

\begin{verbatim}
uv run alembic upgrade head
uv run pytest --cov=abridgeai --cov-report=xml \
  --cov-report=term-missing --cov-fail-under=75
\end{verbatim}

The CI environment provides PostgreSQL and Redis service containers. Its quality gates are linting, security scanning, and the test job; the test job does not claim to execute paid or credentialed voice infrastructure. The configured default marker expression excludes \texttt{destructive} and \texttt{voice\_live} tests.

A separate opt-in harness, \path{tests/integration/test\_voice\_live\_harness.py}, drives English and Vietnamese LiveKit voice scenarios. It requires \texttt{RUN\_LIVEKIT\_VOICE\_TESTS=1}, LiveKit credentials, and an LLM endpoint credential; it checks audible agent output, persisted turns, and cleanup. Because this harness consumes external real-time services, it is intentionally excluded from the ordinary CI suite. This distinction ensures that deterministic state and persistence behavior is continuously verified while live STT/TTS/WebRTC validation remains an explicit operational test.
